% CV master template. Compile with pdfLaTeX (Overleaf or TeX Live).
% Everything below is an example: Alex Chen, Northfield University and all papers and numbers are invented.
% Replace the content with your own. Keep the structure that tools/tailor_cv.py relies on (listed at the top of that script).
\documentclass[a4paper,10pt]{article}

% ---- fonts ----------------------------------------------------------
\usepackage[T1]{fontenc}
\usepackage[sfdefault]{inter}
\usepackage{microtype}
\input{glyphtounicode}\pdfgentounicode=1   % clean text extraction for ATS
\newcommand{\tracked}[1]{\textls[110]{#1}}
\usepackage[hmargin=2.05cm,vmargin=1.7cm,footskip=0.8cm]{geometry}
\usepackage{enumitem}
\usepackage{titlesec}
\usepackage{xcolor}
\usepackage{tikz}
\usepackage{fancyhdr}
\usepackage{lastpage}
\usepackage{needspace}
\usepackage{hyperref}

% ---- palette ------------------------------------------------------
\definecolor{Accent}{HTML}{5B5BD6}
\definecolor{AccentDark}{HTML}{4D4DC4}
\definecolor{Soft}{HTML}{E3E3FD}
\definecolor{Ink}{HTML}{1F2937}
\definecolor{Muted}{HTML}{5D6B82}
\definecolor{Border}{HTML}{E5EAF2}

\hypersetup{colorlinks=true, urlcolor=Accent, linkcolor=Accent, citecolor=Accent,
            pdftitle={Alex Chen - Curriculum Vitae}, pdfauthor={Alex Chen}}

\color{Ink}
\setlength{\parindent}{0pt}
\setlength{\parskip}{0pt}
\raggedright
\hyphenpenalty=10000 \exhyphenpenalty=10000
\widowpenalty=10000 \clubpenalty=10000
\linespread{1.08}

% ---- footer ---------------------------------------------------------
\pagestyle{fancy}
\fancyhf{}
\renewcommand{\headrulewidth}{0pt}
\fancyfoot[R]{\footnotesize\color{Muted}Alex Chen\sep page \thepage\ of \pageref*{LastPage}}

% ---- section heading: accent dot + tracked caps + hairline ----------
\newcommand{\sechead}[1]{%
  \makebox[\textwidth][l]{%
    \tikz[baseline=-0.55ex]{\fill[Accent] (0,0) circle (2.3pt);}\hspace{6pt}%
    {\small\bfseries\color{Accent}\tracked{\MakeUppercase{#1}}}\hspace{10pt}%
    \color{Border}\leaders\hrule height 0.7pt\hfill}}
\titleformat{\section}{\normalfont}{}{0pt}{\sechead}
\titlespacing*{\section}{0pt}{2.6ex}{1.2ex}
\newcommand{\sectionbreak}{\Needspace{9\baselineskip}}

% ---- helpers --------------------------------------------------------
\makeatletter
% dot separator: spacing built in, never breaks before the dot
\newcommand{\sep}{\unskip\nobreak\hspace{0.45em}\textcolor{Muted}{\textperiodcentered}\hspace{0.45em}\ignorespaces}
\newcommand{\muted}[1]{\textcolor{Muted}{#1}}
\newcommand{\pill}[1]{\tikz[baseline=(p.base)]\node[fill=Soft, text=AccentDark, rounded corners=6pt,
  inner xsep=6pt, inner ysep=2.6pt, outer sep=0pt, font=\footnotesize\bfseries] (p) {#1};\hspace{4pt}}
\newcommand{\links}[1]{{\small #1}}   % link row after a publication (hyperref sets the colour)
\newcommand{\me}{\textbf{Alex Chen}}

% left label / right body (no extra gap when it is the first item after a heading)
\newcommand{\entry}[2]{%
  \par\if@nobreak\else\vspace{4pt}\fi
  \begin{minipage}[t]{0.155\textwidth}\raggedright\textcolor{Muted}{\small #1}\end{minipage}%
  \hfill
  \begin{minipage}[t]{0.825\textwidth}\raggedright #2\end{minipage}\par}

% project header: title, subtitle, venue on the right
\newcommand{\project}[3]{\par\Needspace{6\baselineskip}\if@nobreak\else\vspace{5pt}\fi\textbf{#1}\ \muted{#2}\hfill\muted{\small #3}\par\nopagebreak\vspace{1pt}}

\newlist{tightitem}{itemize}{1}
\newlist{publist}{enumerate}{1}
\setlist[publist]{leftmargin=2.3em, labelsep=0.6em, topsep=1pt, itemsep=2.5pt, parsep=0pt, label=\pubnum}
\DeclareRobustCommand{\pubnum}{\textcolor{Muted}{[\the\numexpr\pubtotal+1-\csname c@publisti\endcsname\relax]}}
\newcommand{\pubtotal}{4}  % numbering counts down, newest first
\setlist[tightitem]{leftmargin=1.1em, topsep=1pt, itemsep=1.5pt, parsep=0pt,
  label=\textcolor{Accent}{\raisebox{0.15ex}{\scriptsize\textbullet}}}
\makeatother

%=====================================================================
\begin{document}

% ============================ HEADER ===============================
{\Huge\bfseries Alex Chen}\par
\vspace{2pt}
{\large\color{Muted}PhD Candidate in Computer Science \sep Northfield University}\par
\vspace{5pt}
{\small
Utrecht, the Netherlands \sep +31 6 0000 0000 \sep \href{mailto:alex.chen@example.com}{alex.chen@example.com}\par
\vspace{1.5pt}
\href{https://alexchen.example.com/}{alexchen.example.com} \sep
\href{https://github.com/example}{GitHub} \sep
\href{https://www.linkedin.com/in/example/}{LinkedIn}\par}

% ============================ SUMMARY ==============================
\section{About Me}
I work on time series forecasting for energy systems, where sensors fail, demand patterns shift and models trained last year stop working. I build forecasting and monitoring models that stay accurate when the data changes, and I release the code and benchmarks behind them.

\vspace{4pt}
Open to Machine Learning Engineer and Data Scientist roles, and postdoc positions.\par
\vspace{6pt}
\pill{Time series forecasting}\pill{Anomaly detection}\pill{Distribution shift}\pill{Energy data}\pill{MLOps}\par

% ============================ EDUCATION ============================
\section{Education}
\entry{2022 -- present}{%
  \textbf{PhD in Computer Science}, Northfield University. Defense planned for 2027.\par
  Thesis: \textit{Forecasting Under Shift: Robust Models for Energy Time Series}.\par
  Supervisor Prof. Dr. Jane Doe.}
\entry{2019 -- 2022}{%
  \textbf{MSc in Data Science}, Lakeside Institute of Technology.}

% ============================ EXPERIENCE ===========================
% Role lines group the projects by where they were done.
% tailor_cv.py reorders only the first group: the \project blocks before the \par\Needspace{8\baselineskip} line.
\section{Experience}
{\color{Accent}\bfseries PhD Candidate}\sep Data Systems Lab, Northfield University\hfill\muted{\small 2022 -- present}\par\nopagebreak
\vspace{1pt}
Two first-author projects with released code, on load data from about 12,000 smart meters.\par\nopagebreak

\project{GridCast}{Load forecasting under distribution shift}{Example Conference 2025}
\begin{tightitem}
  \item Forecasting models are trained on past seasons and then run through weather and tariff changes they have never seen. I designed GridCast, which separates slow seasonal structure from fast local changes and retrains only the fast part.
  \item Error at a one-year gap drops by \textbf{12\%} against the strongest baseline, with a public benchmark.
\end{tightitem}

\project{DriftWatch}{Monitoring forecast quality without labels}{under review}
\begin{tightitem}
  \item Ground truth for a forecast arrives days later, so a broken model can run unnoticed. I built a monitor that estimates forecast error from input statistics alone and raises an alert when it leaves the training range.
  \item It flags \textbf{9 of 10} injected sensor faults within one day on the GridCast benchmark.
\end{tightitem}

\project{RiverBench}{A benchmark for hydrological forecasting}{Example Workshop 2024}
\begin{tightitem}
  \item With Sam Lee (lead), we collected ten years of river gauge data from 40 stations and defined forecasting tasks with fixed splits.
\end{tightitem}

\par\Needspace{8\baselineskip}\vspace{8pt}
{\color{Accent}\bfseries MSc Student}\sep Lakeside Institute of Technology\hfill\muted{\small 2019 -- 2022}\par\nopagebreak
\project{MeterClean}{Cleaning smart meter data}{Example Journal 2022}
\begin{tightitem}
  \item I built a pipeline that detects and repairs gaps and spikes in smart meter readings before they reach the forecasting model.
\end{tightitem}

% ============================ SKILLS ===============================
\section{Skills}
\entry{Programming}{Python, SQL, Bash, Git, Linux.}
\entry{ML tools}{PyTorch, scikit-learn, NumPy, Pandas.}
\entry{Methods}{Time series forecasting, anomaly detection, distribution shift, probabilistic models.}
\entry{Languages}{English (fluent), Dutch (basic).}

% ========================== PUBLICATIONS ===========================
\section{Publications}
\begin{publist}
  \item GridCast: Load Forecasting Under Distribution Shift. \me, Jane Doe. \textit{Example Conference 2025}. \links{\href{https://example.com/gridcast}{Paper} \sep \href{https://github.com/example/gridcast}{Code}}

  \item DriftWatch: Monitoring Forecast Quality Without Labels. \me, Jane Doe. Under review.

  \item RiverBench: A Benchmark for Hydrological Forecasting. Sam Lee, \me, Jane Doe. \textit{Example Workshop 2024}.

  \item MeterClean: Cleaning Smart Meter Data at Scale. \me, John Roe. \textit{Example Journal 2022}.
\end{publist}

\end{document}
